\section{Fast simultaneous synthetic butterfly layers}
\label{sec:fast-butterfly-layers}

We now apply the complex network to a common butterfly layer on several
axes. The array has complete address set
\[
 [P]\times[2^K]^D\times[S],\qquad 1\le D\le d,
\]
with $P,S\ge1$, in lexicographic order. Its $D$ consecutive $K$-bit
chunks each select the bit at one common offset $0\le\rho<K$, counted
from the least significant
end. Each address carries a polynomial record with $r$ complex
coefficients, initially in the unit disk on the grid
$2^{-p}\mathbb Z[i]$, with $\Theta(p)$ bits per real or imaginary
component. The intended operation is the tensor product
of $H_0(u,v)=((u+v)/2,(u-v)/2)$ on the $D$ selected bits, acting
coefficientwise and preserving all other address fields.

The parameters will satisfy
\[
 d=\Theta(p^\epsilon),\qquad K=\lfloor d^c\rfloor,
 \qquad r=2^{\Theta(p/d)},\qquad c>0,\quad 0<\epsilon<1.
\]
The exponents and comparison constants are fixed, and we impose their
further restrictions in Proposition~\ref{prop:simultaneous-layer}.
Two consequences explain the layout. The polynomial length $r$ grows
faster than every fixed power of $p$, so a scan of a record pays for
polynomial work on its address. Also $K/\log p\to\infty$, which will
make the exceptional set in a packed change of address bits small enough
to repair by sorting. We will obtain cost $O(Vd^{\lambda'})$, with
$\lambda'<1$ and $V$ the logical bit volume, whereas separate butterfly
factors would require $D$ whole-array passes.

We first construct tensor products of the auxiliary two-point kernel
$C$ from Proposition~\ref{prop:complex-motif-interface}. Its phase
frames express a $C$ tensor product using smaller groups of
selected bits, separated by binary changes of basis. We implement those
changes by moving whole address chunks. Splitting complete rows among
the network's wire roles then reduces the volume of each recursive
call. Finally, we bound the intermediate dyadic values and convert the
completed $C$ layer to the normalized $H_0$ layer. The entire layer is
evaluated exactly before one final truncation.

\subsection{Phase frames for a simultaneous butterfly layer}
\label{subsec:phase-frames}

Hold the unselected address bits fixed and regard the stored complex
values on one wire role as a function of the selected bits.  The phase
identities below express an edge's residual operator as a product of
two-point kernels on selected directions.

Put
\[
 H_0=\frac12\begin{pmatrix}1&1\\1&-1\end{pmatrix},\qquad
 H=\sqrt2 H_0,\qquad \mathsf S=\operatorname{diag}(1,i),\qquad
 Z=\operatorname{diag}(1,-1).
\]
For the bit-flip matrix $X=\bigl(\begin{smallmatrix}0&1\\1&0\end{smallmatrix}\bigr)$,
define
\begin{equation}\label{eq:phase-C}
 C=H\mathsf SH=aI+bX,\qquad a=\frac{1+i}{2},\quad b=\frac{1-i}{2}.
\end{equation}
Direct multiplication gives
\begin{equation}\label{eq:phase-basic-identities}
 H_0=b\mathsf SC\mathsf S,\qquad C^2=X,\qquad C^4=I,\qquad
 C^{-1}=bI+aX=-iZCZ.
\end{equation}
In particular, $C=iZC^{-1}Z$.  These identities use only Gaussian dyadic
coefficients.  The Hadamard matrices in the definition of
$\mathcal C_U$ are used to prove identities; they need not be executed
by the machine.

Consider an edge from label $U_-$ to label $U_+$, and choose an
orthonormal basis $v_1,\ldots,v_{\rho_{\rm edge}}$ of its residual.
Proposition~\ref{prop:complex-motif-interface} gives the frame change
\[
 \mathcal C_{U_+}\mathcal C_{U_-}^{-1}
 =\prod_{j=1}^{\rho_{\rm edge}}(aI+bX_{v_j})^{\epsilon_j},
 \qquad \epsilon_j\in\{1,-1\}.
\]
Here $(X_vf)(x)=f(x+v)$, and we write $C_v=aI+bX_v$.
The number of directional kernels is precisely
$\rho_{\rm edge}$, including on decreasing edges. The next step puts those directions
into coordinate positions without performing a general transpose.

\paragraph{Implementing an arbitrary residual basis.}
For $M\in\operatorname{GL}(m,\mathbb F_2)$, let $\mathcal P_M$ be the
address permutation that moves the value at $x$ to $Mx$.  Thus
$(\mathcal P_Mf)(x)=f(M^{-1}x)$, and
\begin{equation}\label{eq:phase-conjugation}
 \mathcal P_MX_v\mathcal P_M^{-1}=X_{Mv}.
\end{equation}
Extend $v_1,\ldots,v_{\rho_{\rm edge}}$ to a basis of $\mathbb F_2^m$ and take these vectors as
the first $\rho_{\rm edge}$ columns of $M$.  The displayed edge identity
then becomes
\begin{equation}\label{eq:phase-edge-coordinate}
 \mathcal C_{U_+}\mathcal C_{U_-}^{-1}
 =\mathcal P_M\left(\prod_{j=1}^{\rho_{\rm edge}}C_{e_j}^{\epsilon_j}\right)
   \mathcal P_M^{-1},
\end{equation}
where $e_j$ is the $j$th standard basis vector.  The matrix $M$ need not
preserve the dot product: equation~\eqref{eq:phase-conjugation} only
uses its invertibility.

For the array layout used below, take a group of $e=mf$ consecutive
axis chunks and divide it into $m$ consecutive slots of $f$ chunks,
numbered in physical order.
Number the chunks within each slot by $j=0,\ldots,f-1$, starting at
its least significant end.  Write $x_{h,j}$ for the selected bit in
chunk $j$ of slot $h$, where $1\le h\le m$.  Column $j$ is the vector
$(x_{1,j},\ldots,x_{m,j})$.  A frame on this group is the tensor
product of the same phase frame
$\mathcal C_U$ of Proposition~\ref{prop:complex-motif-interface}
on these $f$ columns.

To implement \eqref{eq:phase-edge-coordinate} on the group, first apply
$\mathcal P_M^{-1}$ separately to each column.  For each
$1\le h\le\rho_{\rm edge}$, apply $(C^{\epsilon_h})^{\otimes f}$ to the
$f$ bits $x_{h,0},\ldots,x_{h,f-1}$.  This is one child operation
on the $f$ axes in slot $h$.  Finally apply $\mathcal P_M$ to every
column.  A negative directional kernel still requires only one
forward child, because
\begin{equation}\label{eq:phase-inverse-child}
 (C^{-1})^{\otimes f}
 =(-i)^f Z^{\otimes f}C^{\otimes f}Z^{\otimes f}.
\end{equation}
The rightmost phase is applied first.  All inverse wrappers are address
signs and constant fourth-root phases, so they add no recursive calls.

Gaussian elimination expresses any $M\in\operatorname{GL}(m,\mathbb F_2)$
as a product of at most $m(m-1)+3m$ elementary row additions: one clears
at most $m-1$ other entries per pivot, and each row swap is three row
additions.  Reversing the elimination gives the required product since
each row addition is an involution.  The number is a constant because
$m$ is fixed.  Adding row $h'$ to row $h$ in every column applies
$x_{h,j}\leftarrow x_{h,j}\mathbin\oplus x_{h',j}$ for every
$0\le j<f$.  Thus one row addition changes corresponding selected
positions in two whole slots.  The packed address operation proved
below performs all of these XORs together.

\paragraph{The exact recursive contract.}
Apply the finite network with these column frames to a group of
$e=mf$ selected axes.  Summing its residual dimensions gives exactly
$s_{\rm c}$ child calls, each on the $f=e/m$ axes of one slot.
The source and sink $Z_{a,f}$ signs, together with the scalar
$i^{27f}$ at each affected sink, turn every endpoint kernel into a
forward one.  Negating the prescribed physical bank
and undoing its exchange returns each result to its original role.
Proposition~\ref{prop:complex-motif-interface} therefore gives
$C^{\otimes e}$ on every role, including roles whose initial scratch
values were arbitrary.  Formula~\eqref{eq:phase-inverse-child} supplies
each inverse child from one forward call, and none of the corrections
adds a child.  The remaining operations are a fixed number of binary
row additions and coefficientwise scalar operations.

\paragraph{Returning to normalized butterflies.}
For an address $x$ in a full layer, write $x_1,\ldots,x_D$ for its
selected bits.
The diagonal operator $\mathcal S_D$ multiplies the value stored at
an address by the phase determined by those bits:
\[
 (\mathcal S_D z)(x)=i^{x_1+\cdots+x_D}z(x).
\]
Here the unselected address bits are unchanged, and the same phase is
applied to every coefficient stored at that address.  Applying
\eqref{eq:phase-basic-identities} on all selected coordinates gives
\begin{equation}\label{eq:phase-butterfly-layer}
 H_0^{\otimes D}=b^D\mathcal S_D C^{\otimes D}\mathcal S_D.
\end{equation}
The scalar is
\[
 b^D=(-i/2)^{\lfloor D/2\rfloor}b^{D\bmod2},
\]
so it is a binary shift and fourth-root phase, with at most one further
factor $b$.  All phases, inverse wrappers, and scalar corrections above
are exact Gaussian dyadic operations.  Thus, once the physical binary
basis changes and recursive stream layout are supplied, the complete
parallel butterfly layer is exact up to its prescribed final truncation.

\subsection{Packed changes of selected address bits}
\label{sec:packed-selected-bits}

We implement a binary row addition on corresponding selected bits in two
address chunks. A third chunk supplies temporary address bits; its arbitrary
initial value must be restored. The construction first uses a constant number
of cyclic shifts of whole chunk ranges, then corrects an explicitly bounded set of
addresses. The correction depends only on addresses, so the time bound is
worst-case over all payloads.

We use three previously established streaming operations.
Lemma~\ref{lem:chunk-swap} exchanges two disjoint address chunks of
equal width $L$ in $O(VL^\tau)$ time, where $V$ is the current stream
volume and $0<\tau<1$ is fixed.  The construction in
Lemma~\ref{lem:ordered-affine-streams} cyclically shifts the complete
range of a target chunk by an offset determined by preceding fields in
$O(V)$ time apart from calculating the offsets.  Finally,
Lemma~\ref{lem:elementary-streams} applies any fixed permutation to the
two-bit value at two specified address positions in $O(V)$ time.
Each operation uses a fixed number of one-dimensional
tapes and restores the order of all spectator fields.

The asymptotic simplification in the following lemma describes a fixed
family of parameters as \(p\to\infty\), after the exponents and comparison
constants have been fixed.  The statement that \(R\) exceeds every fixed
polynomial in \(p\) means that \(R/p^C\to\infty\) for each fixed \(C\).
Implied constants and eventual cutoffs may depend on the fixed choices
and on when the family's asserted bounds begin to hold.  In particular,
each explicit bound is uniform over all of its declared inputs;
its asymptotic simplification is uniform over the allowed row counts
and payloads within a fixed family.

\begin{lemma}[Packed selected-bit addition on one rectangle]
\label{lem:packed-selected-bit-rectangle}
Let $R\ge1$, $K\ge6$, $f\ge1$, $L=fK$, and $0\le\rho<K$.
Let $N_0,N_1,N_2,N_3$ be positive integers, and let
$(\chi_1,\chi_2,\chi_3)$ be a permutation of $(x,y,z)$.
Consider $R$-bit records in lexicographic order on
\[
 [N_0]\times[2^L]_{\chi_1}\times[N_1]\times[2^L]_{\chi_2}
 \times[N_2]\times[2^L]_{\chi_3}\times[N_3].
\]
The subscript names the chunk at that position. The input supplies
$N_0,\ldots,N_3,R,K,f,\rho$ as a fixed number of canonical binary
integers, together with the finite permutation of the chunk names, as
in Section~\ref{sec:streams}. Put
\[
 M=2^{3L}\prod_{h=0}^3N_h,\qquad V=MR,\qquad
 A_{\rm rect}=\lceil\log_2(2V)\rceil.
\]
At the selected positions $j_i=\rho+iK$, numbered from the least
significant bit, one fixed multitape procedure performs
\[
 y_{j_i}\longleftarrow y_{j_i}\mathbin\oplus x_{j_i}
 \quad(0\le i<f).
\]
It preserves the spectator coordinates, all other chunk bits, and every
record payload; in particular, it restores the entire chunk $z$ for
every initial value.
With $\delta_{\rm bad}=\min\{1,80(f-1)2^{-K}\}$, its time is
\begin{equation}
\label{eq:packed-rectangle-time}
 O\!\left(V(L^\tau+1)+MA_{\rm rect}^3+
       \delta_{\rm bad}MA_{\rm rect}(R+A_{\rm rect})\right).
\end{equation}
When $A_{\rm rect}=O(p)$, $R$ exceeds every fixed polynomial in $p$,
$f\le d\le p$, and $K/\log p\to\infty$, this is
$O(V(L^\tau+1))$, uniformly over the permitted lengths and payloads
in the fixed family.
\end{lemma}

\begin{proof}
We first construct eight rotations that have the required effect away
from a small set of addresses. We then correct exactly that set.
For a Boolean control $c$ and arbitrary integers $u,w$, perform the following
eight operations in the displayed chronological order. Every predicate
uses the current values of the integers.
\begin{equation}
\label{eq:eight-packed-additions}
\begin{array}{rcl@{\qquad}l}
u&\leftarrow&u-c[\,w\equiv0\pmod2\,],&(1)\\
w&\leftarrow&w-c[\,u\equiv0\pmod2\,],&(2)\\
u&\leftarrow&u+c[\,w\equiv0\pmod2\,],&(3)\\
w&\leftarrow&w+c[\,u\equiv0\pmod2\,],&(4)\\
w&\leftarrow&w+c[\,u\equiv1\pmod2\,],&(5)\\
u&\leftarrow&u+c[\,w\equiv1\pmod2\,],&(6)\\
w&\leftarrow&w-c[\,u\equiv1\pmod2\,],&(7)\\
u&\leftarrow&u-c[\,w\equiv1\pmod2\,].&(8)
\end{array}
\end{equation}
If $c=0$, these operations do nothing. If $c=1$, the first four operations
induce the cycle $(00\ 01\ 10)$ in every square
\[
 \{2a,2a+1\}\times\{2b,2b+1\},\qquad a,b\in\mathbb Z,
\]
where the corner labels record the parities of $(u,w)$. The last four
induce $(01\ 10\ 11)$. Applying the first cycle and then the second gives
$(00\ 10)(01\ 11)$. These facts follow by substituting the four parity
pairs into the displayed operations; translation by $(2a,2b)$ does not
change any predicate. Consequently the complete output is
\begin{equation}
\label{eq:eight-packed-result}
 (u',w')=
 \bigl(u+c(1-2(u\bmod2)),\ w\bigr).
\end{equation}
Here $u\bmod2\in\{0,1\}$, also for negative $u$. Thus the parity of $u$
is flipped precisely when $c=1$, both quotients under division by two are
preserved, and the arbitrary auxiliary integer $w$ is restored. Each
coordinate is updated four times by at most one, so its displacement
from its initial value is at most four, even if an intermediate point
leaves its starting quotient square. The weaker bound eight will suffice below.

Omit the most significant selected position $j_{f-1}$ temporarily. For
each $0\le i<f-1$, the bits in positions
$j_i,\ldots,j_i+K-1$ form a $K$-bit segment. When $f>1$, these segments
are disjoint, and the highest one ends immediately below $j_{f-1}$.
Write their initial values in $y$ and $z$ as
\[
 u_i=2g_i+a_i,\qquad w_i=2h_i+b_i,
 \qquad a_i,b_i\in\{0,1\}.
\]
The values $g_i,h_i\in\{0,\ldots,2^{K-1}-1\}$ are the guard values. The
control for this pair of segments is the selected bit $c_i=x_{j_i}$.

Apply each line of \eqref{eq:eight-packed-additions} simultaneously to all
these segment pairs. Physically, a line updating $y$ by sign $\varepsilon$
and testing parity $b$ of $z$ is the single cyclic shift
\begin{equation}
\label{eq:packed-cyclic-shift}
 y\longleftarrow y+
 \varepsilon\sum_{i=0}^{f-2}2^{j_i}
 [\,x_{j_i}=1\,][\,z_{j_i}=b\,]
 \pmod{2^L}.
\end{equation}
For a line updating $z$, interchange $y$ and $z$ in this formula.
The shift offset never depends on its target chunk, so each line is a
bijection of the rectangle. Let $S_0$ be their composition. It fixes
the spectator coordinates.

Let $\mathcal B$ consist of the addresses for which
at least one used guard value, in either $y$ or $z$, fails
\begin{equation}
\label{eq:packed-guard-interval}
 10\le g\le 2^{K-1}-11.
\end{equation}
At an address outside $\mathcal B$, each segment initially lies in
$[20,2^K-21]$. Even displacement eight keeps it inside
$[0,2^K-1]$. Inductively, each packed addition therefore agrees with all
the independent segment additions: no segment carries into or borrows
from the next segment. When $f>1$, the bits below $j_0$ and above the
highest used segment are untouched as well. By
\eqref{eq:eight-packed-result}, $S_0$
then agrees with the ideal permutation $T_0$ that performs the required
XORs at $j_0,\ldots,j_{f-2}$ and fixes all other address coordinates.

There are exactly twenty excluded values for each guard in
\eqref{eq:packed-guard-interval}. Since its complete range has
$2^{K-1}$ values, the union bound over the $2(f-1)$ guards gives
\begin{equation}
\label{eq:packed-bad-fraction}
 |\mathcal B|\le\delta_{\rm bad}M.
\end{equation}
Indeed, the complete $y,z$ ranges give the same guard distribution at
each choice of spectator coordinates. No record value enters this count.

The ideal map $T_0$ preserves all guard values and therefore preserves
$\mathcal B$. Both $T_0$ and $S_0$ are bijections, and they agree off
$\mathcal B$. Taking complements in the finite rectangle gives
\[
 S_0(\mathcal B^c)=T_0(\mathcal B^c)=\mathcal B^c,
 \qquad S_0(\mathcal B)=\mathcal B.
\]
Thus $T_0S_0^{-1}$ fixes $\mathcal B^c$ and permutes $\mathcal B$.
It remains to implement the rotations and this correction on tapes.

For one rotation, move its target to the rightmost of the three chunk
positions if necessary. Let $P^{\rm sw}$ be the product of the record
field lengths before the earlier exchanged chunk, let $G^{\rm sw}$
be the product between the exchanged chunks, and let $B^{\rm sw}$ be
$R$ times the product after the later one. Empty products are one.
The bit array has exactly the shape
\[
 [P^{\rm sw}]\times[2^L]\times[G^{\rm sw}]
       \times[2^L]\times[B^{\rm sw}].
\]
The third chunk lies inside one collapsed spectator interval.
Pass only $L$ and these three positive canonical lengths to
Lemma~\ref{lem:chunk-swap}. Each length is at most $V$, so the local
descriptor has $O(A_{\rm rect})$ bits. The swap costs $O(VL^\tau)$, including
its own padding and cleanup.

The equal chunk lengths leave the field-length list unchanged by a
completed swap. Keep the permutation of the logical names $x,y,z$ in
finite local state. With the target now after its two controls, collapse
the remaining intervals to obtain
\[
 [\Lambda_0]\times[2^L]_{\rm control\,1}\times[\Lambda_1]
 \times[2^L]_{\rm control\,2}\times[\Lambda_2]
 \times[2^L]_{\rm target}\times[B_{\rm rot}],
\]
where $B_{\rm rot}$ includes the payload and every displayed spectator
length is positive. Their product with the chunk lengths is $V$.
Together with $K,f,\rho$ and the finite name permutation, these fixed
canonical lengths form an $O(A_{\rm rect})$-bit local descriptor. The offset
routine receives the two current $L$-bit logical control values. A scan
of at most $f\le L$ selected positions evaluates
\eqref{eq:packed-cyclic-shift} in $L^{O(1)}$ time, without reading the
target or spectator values. Lemma~\ref{lem:ordered-affine-streams}
therefore gives an $O(V)$ rotation: its offset work is absorbed by the
$2^L$ nonempty target blocks in each fiber. Undo the swap after the
rotation. The local descriptors and name permutation are updated at
completed swap boundaries; the primitives use their own descriptors
inside padded layouts.
There are at most sixteen swaps and eight rotations in $S_0$.

To implement the correction, scan the stream with a record-rank counter.
All chunk names are back in their original positions. Dividing the
current rank through the fixed seven field lengths recovers the three
chunk values and their mixed-radix record strides. The strides exclude
the $R$ payload bits. All involved integers have $O(A_{\rm rect})$ bits, so
schoolbook arithmetic, descriptor copying, and resets take at most
$O(A_{\rm rect}^3)$ per record. Test membership in $\mathcal B$ and extract just
those records. For a current exceptional address $q$, compute
$q'=T_0(S_0^{-1}(q))$ by reversing the eight modular additions with
negated offsets, then applying the selected XORs directly. Each reversed
mask uses the current controls of that reversed line. There are at most
$f\le L\le A_{\rm rect}$ selected positions, so this calculation also fits
$O(A_{\rm rect}^3)$ time.

If $q_h,q'_h$ are the current and destination values of chunk
$h\in\{x,y,z\}$ and $\sigma_h$ is its record stride, the fixed
spectator coordinates give
\[
 \operatorname{rank}(q')=\operatorname{rank}(q)+
       \sum_{h\in\{x,y,z\}}(q'_h-q_h)\sigma_h.
\]
The differences are signed. Attach this destination rank, padded to
$\lceil\log_2M\rceil\le A_{\rm rect}$ bits, to the extracted record, and mark its
hole on a fixed track of the full stream. Hold the full stream untouched
while stable binary radix sorting scans only the extracted records and
keys, least significant key bit first. Since $T_0S_0^{-1}$ permutes
$\mathcal B$, these keys are distinct and are exactly the ranks of the
marked holes. One forward full-stream scan reinserts the sorted records.
This implements $T_0$ everywhere.

Finally apply the two-bit XOR to $x_{j_{f-1}},y_{j_{f-1}}$ by
Lemma~\ref{lem:elementary-streams}, splitting those chunks at the
selected positions and collapsing the other intervals. Its descriptor
again has a fixed number of positive lengths, each at most $V$.
This costs $O(V)$ and supplies the omitted selected position. When
$f=1$, $S_0,T_0$ are identities, $\mathcal B$ is empty, and this last
step is the entire algorithm.

The movement costs $O(V(L^\tau+1))$. A fixed number of full record scans
costs $O(MA_{\rm rect}^3+V)$, and the at most $A_{\rm rect}$ radix passes cost
$O(|\mathcal B|A_{\rm rect}(R+A_{\rm rect}))$; skip them when the extracted bank is empty.
The initial descriptor has a fixed number of canonical integers of
$O(A_{\rm rect})$ bits, so its preparation is covered by $O(MA_{\rm rect}^3)$.
All counters, positioning, rewinds, and cleanup are included.
Equation~\eqref{eq:packed-bad-fraction} proves
\eqref{eq:packed-rectangle-time} on a fixed number of tapes.
After division by $V=MR$, its additional terms are bounded by
$A_{\rm rect}^3/R$ and
$80(f-1)2^{-K}A_{\rm rect}(1+A_{\rm rect}/R)$. Under the stated family bounds
the first tends to zero, and the second is
$O(p^2 2^{-K}(1+A_{\rm rect}/R))=o(1)$. This proves the final assertion.
\end{proof}

\paragraph{Application to a fixed binary change of basis.}
An invertible matrix over $\mathbb F_2$ of fixed order $m\ge3$ is a
product of a fixed number of elementary row additions; a row exchange
is three such additions.  In one invocation of the finite network on
$e=mf$ axes, slot $h$ consists of $f$ consecutive $K$-bit axis chunks.
It therefore has address-bit width $L=fK$, and its selected bits
$x_{h,j}$ occupy positions $\rho+jK$ for $0\le j<f$.
To add row $h'$ to row $h$, use slot $h'$ as $x$,
slot $h$ as $y$, and any third, distinct slot as $z$ in
Lemma~\ref{lem:packed-selected-bit-rectangle}.  The lemma performs
$x_{h,j}\leftarrow x_{h,j}\mathbin\oplus x_{h',j}$ for every $j$ in
one operation and restores the third slot, whatever its initial
address value.

Under the asymptotic hypotheses in the final assertion of that lemma,
the complete change of basis and its inverse therefore cost
$O(V((eK)^\tau+1))$ on a stream of volume $V$.  Here $e$ counts
axis chunks, while $eK$ counts address bits; the fixed factor $m$
between $eK$ and the slot width does not affect the bound.  The same
conclusion holds after the row padding and splitting below, because
every role stream retains a complete suffix with the same fixed field
lengths in every row.

\subsection{Batching a common layer on a fixed number of tapes}

We now turn the finite complex network into an algorithm for a whole
layer.  Throughout this subsection its constants are denoted by
\(m,W,s\): the network has \(W\) wires, its binary label space has
dimension \(m\), and the sum of the dimensions of its edge residuals is
\(s<Wm\).  These are the constants of the complex network; they need not
equal the constants used to prove the chunk-swap lemma.

Use the complete layout and parameter family specified at the start of
the section, with fixed rational exponents \(c,\epsilon\) to satisfy
the recurrence and guard conditions below. Store each polynomial in
coefficient order with a common width \(w=\Theta(p)\) per real or
imaginary component, and assume the complete address and its binary
shape descriptor have length \(O(p)\). Thus an input record has
\(2rw\) bits. Write \(V\) for the current logical bit volume, excluding
temporary work streams. We first count the work on \(C^{\otimes D}\)
assuming one common exact coefficient format of width \(O(p)\); the
guard calculation below then supplies that format.

\paragraph{An individual selected-bit kernel.}
Write each input coefficient once in this common exact format and retain it until
the final truncation.  The input width is \(\Theta(p)\), so the
initial copy costs \(O(V)\) and the guarded stream has volume \(O(V)\).

To apply one \(C\) kernel, use the split from
Lemma~\ref{lem:elementary-streams} at the selected address bit.
Call the two streams \(0\) and \(1\) according to that bit.
Because that bit has its complete two-element range, the two streams
have the same remaining address set in the same order.  Their aligned
record addresses therefore differ only in the selected bit, and their
coefficients are still ordered by the same polynomial index.  Scan
each pair of records in that order.  For corresponding complex
coefficients \(u\) from stream \(0\) and \(v\) from stream \(1\), compute
\begin{equation}\label{eq:individual-C-pair}
 C(u,v)=
 \left(\frac{u+v+i(u-v)}2,\,
       \frac{u+v-i(u-v)}2\right).
\end{equation}
Addition, subtraction and halving use the fixed-point conventions of
Section~\ref{sec:streams}; multiplication by \(i\) swaps the real
and imaginary components and changes one sign.  The common guarded
width makes these operations exact.  The coefficient scan costs
\(O(rp)\) per pair of records, linear in their guarded bit length
and including any reversal of bits inside a coefficient.  Write the
first result to stream \(0\) and the second to stream \(1\), both in
coefficient order, then merge
according to the selected bit.  The merge restores the original
lexicographic address order.  The elementary split and merge include
their counters and cleanup and cost \(O(V)\); the paired arithmetic
scan also costs \(O(V)\).  Thus \(e\) individual kernels cost
\(O(Ve)\) on any stream with these complete ranges and record format.

First consider an internal node containing \(e=m^k\) consecutive
active chunks, with \(k\ge1\).  Split it into \(m\) consecutive slots,
each containing \(e/m\) chunks.
The phase construction supplies one recursive call on a slot for each
vector in an edge residual basis, hence exactly \(s\) calls in one
invocation of the network.  If the logical volume entering this node
is \(V\), binary changes of basis, their inverses, scalar gates,
endpoint corrections and reassembly of the physical banks cost
\[
 O\bigl(V((eK)^\tau+1)\bigr)
\]
at this node.  Each slot has \((e/m)K\) address bits, so
Lemma~\ref{lem:packed-selected-bit-rectangle} gives the stated bound for a row addition.  The
network and all of its \(m\)-dimensional residual bases are fixed,
so the total number of row additions is fixed as well.  At the same
guarded width, the scalar gates cost a constant
number of scans of the current volume.  Address and descriptor
work bounded by any fixed polynomial in \(p\) per record is also
absorbed: its ratio to the record's \(\Theta(rp)\) bits tends to zero.

\paragraph{Where the rows come from.}
Put
\[
 q_0=\lceil\log_2 W\rceil\,
       \lceil\log_m(2d)\rceil,
 \qquad k_0=\lceil\log_m d\rceil.
\]
If \(D\le q_0\), use individual \(C\) kernels, at cost \(O(V\log d)\).
Otherwise apply \(C\) individually on the selected bits of the first
\(q_0\) chunks, in \(O(V\log d)\) time.  These chunks remain a
complete consecutive address block before all the remaining chunks;
regard their concatenated bits as one row index.  Its range has
length \(R_{\rm row}=2^{q_0K}\ge W^{k_0}\), since \(K\ge1\) for sufficiently
large \(p\) and \(q_0K\ge k_0\log_2W\).  Within each preceding prefix, add
the same number of zero rows to reach the next multiple of \(W^{k_0}\);
the volume grows
by a factor at most two.  At any recursive depth \(j\le k_0\), the
row count is divisible by \(W^{k_0-j}\).

At a node, write a row index as \(u=Wg+w\), with \(0\le w<W\),
and send that entire row to physical wire role \(w\), indexed there
by \(g\).  Every role stream has the same remaining address shape.
Thus a pointwise gate reads the records with the same \(g\) and the
same remaining address in its touched roles.  The scalar network
allows these role values to be arbitrary, so the split supplies its
scratch roles from existing rows.  Each recursive edge call acts on
one role stream with volume exactly \(1/W\) of its parent's logical
volume.

Splitting the row index removes no address within a row. Every row in
every role stream, including a padded row, retains the complete suffix
\([2^K]^{D-q_0}\times[S]\) in the same field order and with the same
lengths; inactive chunks remain complete spectator fields. The row
counts also remain equal across preceding prefixes, decreasing by
the same factor \(W\) at each split. The network is invoked only on
these equal rows: their common suffix gives the aligned role addresses
and exact \(V/W\) child volumes above. For each packed call, collapse
the preceding prefix and row index into one field, and collapse the
spectator intervals around its three slots. This gives the rectangle of
Lemma~\ref{lem:packed-selected-bit-rectangle}, whose guard count and
deterministic repair apply to every child.

The \(W\) role streams use a fixed number of tapes.  Before entering a
child, push the other role streams onto a LIFO work tape, copy the
active stream to a reserved child input area, and clear or reuse the
role areas.  On return, copy the result back and pop the parked
streams, reversing each popped stream when necessary to restore its
order.  The stack head is at its top at every call boundary.  Pushing,
popping, copying, resetting and clearing at this node cost its own
volume times a fixed constant; no step traverses an ancestor's parked
data.  The call descriptor and program counter use a separate stack.
The number of child calls is the fixed integer \(s\), so these costs
are included in the node overhead above.  The fixed-tape chunk-swap
subroutine has its own fixed work areas, reused on every call.

For each packed call, padding has increased the row count by at most
a factor two and row splitting can only decrease it.  Its active
address still has \(O(p)\) bits. The four collapsed spectator lengths
therefore have \(O(p)\)-bit descriptions. Its record-width parameter
describes a polynomial of \(\Theta(rp)\) bits and has
\(O(\log(rp))=O(p)\) bits as well. These lengths, together with
\(K,f,\rho\) and the finite chunk-name permutation, are precisely the
local descriptor required by
Lemma~\ref{lem:packed-selected-bit-rectangle}. Both the record count
and record width have \(O(p)\)-bit lengths, so its bound uses
\(A_{\rm rect}=\lceil\log_2(2V)\rceil=O(p)\). A selected-position
mask has at most \(DK=O(p)\) bits.

The binary descriptions of all growing shape parameters have
polynomial length in \(p\).  Each nonempty child stream still contains
a complete length-\(r\) polynomial record.  Its volume therefore
absorbs polynomial descriptor setup at that node, as well as
polynomial address work per record.  Even if every node is counted
separately, there are only \(d^{O(1)}\) nodes: the branching factor
\(s\) is fixed and the depth is at most
\(\lceil\log_m d\rceil\).  Any additional fixed power of \(p\)
from this count is still dominated by \(r\).

Every completed network invocation applies \(C^{\otimes e}\)
separately to each input row.  Its selected bits lie in the
remaining chunks, disjoint from the bits used by the earlier
preprocessing.  On the genuine rows, those two completed operations
therefore commute.
Each padded zero row is again zero when an invocation finishes, so the
added rows can be deleted after the invocations on the remaining chunks
have completed. Intermediate scratch values may be arbitrary.

\paragraph{Unrolling the recurrence.}
Fix \(\beta=1/2\), and choose real exponents satisfying
\[
 0<\tau<1,\qquad
 \max\{0,\log_m(s/W)\}\le\sigma<1,
 \qquad \max\{\tau,\sigma\}<\lambda<1,
 \qquad \tau(1+c/\beta)<\lambda.
\]
The exponents can all be chosen rational, with a strict upper bound
in the condition on \(\sigma\).  Stop recursion when
\(e<d^\beta\) and apply the \(e\) \(C\) kernels individually.  If
\(F(e)\) denotes time divided by the current logical stream volume,
then
\[
 F(e)\le (s/W)F(e/m)+O((eK)^\tau+1),
 \qquad F(e)=O(e)\quad(e<d^\beta).
\]
Indeed, the \(s\) child calls are executed on streams of volume
\(V/W\), so their total cost is \(s(V/W)F(e/m)\).  Dividing by
the current volume \(V\) gives the factor \(s/W\).
At an internal node \(e\ge d^\beta\), hence
\(K\le d^c\le e^{c/\beta}\).  The overhead is therefore
\(O(e^{\tau(1+c/\beta)}+1)=O(e^\lambda)\).
Put \(a=s/W\).  Since \(a\le m^\sigma<m^\lambda\), the internal
costs sum to at most
\[
 C e^\lambda\sum_{j\ge0}(a/m^\lambda)^j=O(e^\lambda).
\]
If the leaves after \(j\) levels contain \(u=e/m^j<d^\beta\)
active chunks each, their combined normalized
contribution is at most
\[
 O(a^ju)\le O(e^\sigma u^{1-\sigma})
          \le O(e^\sigma d^{\beta(1-\sigma)}).
\]
The first inequality uses
\(a^j\le m^{j\sigma}=(e/u)^\sigma\).  It also covers \(a<1\)
by taking \(\sigma\ge0\).

Write \(D-q_0=\sum_j a_jm^j\) in base \(m\), with \(0\le a_j<m\).
Partition the remaining chunks into \(a_j\) consecutive pieces
containing \(m^j\) chunks for each \(j\).
There are at most \((m-1)(1+\lfloor\log_m d\rfloor)\) pieces.
Adding their costs and the row preprocessing gives
\[
 O\!\left(\log(2d)
    \bigl[d^\lambda+d^{\sigma+\beta(1-\sigma)}+1\bigr]\right)
   =O(d^{\lambda'})
\]
for any fixed
\[
 \max\{\lambda,\sigma+\beta(1-\sigma)\}<\lambda'<1.
\]
The case \(D\le q_0\), in which preprocessing would consume all
active chunks, was handled by individual kernels.  Thus all constants
and all tape counts are independent of \(p,d,r\) and the input data.

\subsection{An explicit guard-width bound}

The accelerated network is evaluated exactly until an entire
normalized butterfly layer is complete.  Its intermediate operators
need not be contractions.  We now give explicit constants that bound
denominator and magnitude growth within one layer.

For the complex network, let
\[
 E=64(W+m+1)^3,\qquad B=s+E,
 \qquad \nu=\min\{j\ge1:m^j\ge B\},
\]
and set
\[
 C_1=\max\{20,\nu+3\},\qquad C_0=128mB^2,
 \qquad \Delta=\lceil C_0d^{C_1}\rceil.
\]
For the stated \(h=100\) construction, the complex network has
\[
\begin{split}
 m&=1{,}000{,}000,\\
 W&=1{,}873{,}807{,}244{,}643{,}542{,}670{,}000,\\
 s&=1{,}873{,}807{,}244{,}636{,}671{,}267{,}308{,}000{,}000.
\end{split}
\]
These values give \(m^{10}<B<m^{11}\), hence \(\nu=11\) and
\(C_1=20\).  In fact the bound below needs only exponent fourteen;
the choice twenty leaves an explicit additional margin.

We measure arithmetic dependency depth using complex additions,
subtractions, halvings, sign changes and multiplication by \(i\).
For a group of \(e\) axes, let \(A(e)\) bound the number of these
operations on any chain from an input coefficient to an intermediate
coefficient in its recursive \(C\)-layer.  A chain follows an operand
into a later arithmetic result.  Work on other coefficients or records
that do not feed this value lies on other chains, even though the
fixed-tape machine executes that work sequentially.  To obtain a
conservative bound, we concatenate all scalar gates and all recursive
edge calls within one network invocation.

The scalar schedule \eqref{eq:motif-schedule} touches a data wire
at most four times in each of the three stages.  A scratch wire belongs
to one invocation and is touched at most four times there.  Thus every
physical wire occurs in at most twelve scalar gates, and there are
at most \(12W\) gates.  A gate has at most \(W\) inputs and outputs
and coefficients in \(\{0,\pm1,\pm1/2\}\).  One output is a sum of
at most \(W\) terms, each requiring at most one halving and one
addition or subtraction.  Evaluating all outputs from saved inputs
therefore takes at most \(2W^2\) elementary operations.

For an inverse child, each \(Z^{\otimes f}\) is one sign chosen by
the parity of the selected address bits, and the scalar fourth-root
phase takes at most one sign change and one multiplication by \(i\).
This gives at most \(4s\) operations for the child corrections.
The source, sink and bank corrections take at most \(4W+4\) more.
Calculating the address parities and phase exponents is address or descriptor
work, not arithmetic on coefficients.  Since \(s<Wm\), the total
\(24W^3+4s+4W+4\) is less than
\(E=64(W+m+1)^3\).  Binary address permutations add no coefficient
arithmetic at their interfaces.

The paired evaluation in \eqref{eq:individual-C-pair} uses at most
eight elementary operations: form the sum and difference, multiply
the difference by \(i\), then form the two half-sums.  We may therefore use
\[
 A(e)\le sA(e/m)+E
 \quad\hbox{at an internal node},\qquad A(e)\le8e
 \quad\hbox{at a leaf}.
\]
The depth estimate concatenates the \(s\) calls, so it uses \(s\).
The factor \(1/W\) in the earlier time recurrence accounted for
their smaller stream volumes.

There are at most \(k\le\lceil\log_m d\rceil\) internal levels.
For any bound \(x\ge1\), \(sx+E\le(s+E)x=Bx\).
Starting with \(8e\le8d\le8dB\) at a leaf and applying this
inequality through the levels gives the first bound below.
For the second, use \(B\le m^\nu\) and
\[
 B^{k+1}\le B^{\log_m d+2}
          =B^2d^{\log_m B}\le B^2d^\nu.
\]
Thus
\[
 A(e)\le 8dB^{k+1}\le8B^2d^{\nu+1}.
\]
For the sufficiently large sizes under consideration, the earlier
piece count is at most \(md\).  Concatenating their depths contributes
at most \(8mB^2d^{\nu+2}\).  In either preprocessing branch, at most
\(d\) individual kernels are used, contributing at most \(8d\).
Splitting rows, padding, copying and removing padding add no
coefficient arithmetic.

Each of the two outer diagonals applies a fourth-root phase and costs
at most two elementary phase operations.  Put
\(t=\lfloor D/2\rfloor\).  Since \(b^2=-i/2\), the scalar \(b^D\)
is \((-i)^t2^{-t}b^{D\bmod2}\).  Its phase costs at most two
operations, and the possible factor \(b\) is computed as
\((z-iz)/2\) in three operations.  The shift by \(t\) consumes
\(t\) fractional bits, so we charge it \(t\) units of denominator
growth even though one tape scan performs the shift.  Charging these
same units to magnitude growth gives a common upper bound for both:
\[
\begin{aligned}
 8mB^2d^{\nu+2}+8d+t+9
 &\le 8mB^2d^{\nu+2}+18d\\
 &\le 26mB^2d^{\nu+2}\\
 &<128mB^2d^{\nu+3}\le C_0d^{C_1}.
\end{aligned}
\]
Here \(d\ge1\), \(mB^2\ge1\), and \(C_1\ge\nu+3\).

To see how this count controls numerical size, start with disk-valued
coefficients in \(2^{-p}\mathbb Z[i]\).  A value at depth \(\ell\)
lies in \(2^{-(p+\ell)}\mathbb Z[i]\) and has modulus at most
\(2^\ell\).  This follows by induction: addition or subtraction
can at most double the maximum modulus of its operands; halving adds
at most one denominator bit; and signs and multiplication by \(i\)
preserve both bounds.  The charge of \(t\) for the final shift is
the same as \(t\) successive halvings.  Thus reserving \(p+\Delta\)
fractional bits and \(\Delta+3\) integer and sign bits per real
component represents every intermediate coefficient exactly and
prevents overflow.  Shifts consume the reserved fractional zeros
and discard no bits.

An address-swap subroutine may temporarily hold XOR combinations of
the fixed-width encoding bits.  These strings are not used as
coefficients during the swap.  It returns an exact permutation of
the encodings before coefficient arithmetic resumes, so these
internal bit operations cause no numerical rounding or denominator
growth.

Choose \(\epsilon C_1<1\).  Since \(d=\Theta(p^\epsilon)\),
\(\Delta=o(p)\), and the full coefficient width is \(O(p)\), as
required in the time analysis.  Addition, signed shifts, copying and
phase changes consequently take \(O(p)\) bit operations per
coefficient.

\begin{proposition}[Simultaneous normalized butterfly layer]
\label{prop:simultaneous-layer}
Fix
\[
 \tau=\sigma=1-2^{-50},\qquad \beta=\frac12,\qquad C_1=20.
\]
Let $c,\epsilon,\lambda,\lambda'$ be fixed positive rational numbers
satisfying
\[
 \max\{\tau,\sigma\}<\lambda<1,\qquad
 \tau(1+c/\beta)<\lambda,\qquad
 \max\{\lambda,\sigma+\beta(1-\sigma)\}<\lambda'<1,
 \qquad \epsilon C_1<1.
\]
Fix positive constants $a_d\le b_d$ and $a_r\le b_r$, and fix
positive constants $C_M,C_{\rm desc},C_w$ with $C_w>1$.
For these fixed choices
there are a cutoff $p_0$, one fixed finite-alphabet multitape procedure,
and an implicit time constant with the following guarantee. The cutoff
and time constant may depend on the fixed rational parameters
and these constants. For every choice of positive integers $p,d,r$
with $p\ge p_0$ and
\[
 a_dp^\epsilon\le d\le b_dp^\epsilon,\qquad
 2^{a_rp/d}\le r\le 2^{b_rp/d},
\]
put $K=\lfloor d^c\rfloor$, and let $1\le D\le d$, $P,S\ge1$, and
$0\le\rho<K$ be integers. Let $w\ge1$ be an integer component width.
The input supplies on a descriptor tape
the self-delimiting header
\[
 \mathcal E=
 \Gamma(p)\Gamma(d)\Gamma(r)\Gamma(D)\Gamma(P)\Gamma(S)
 \Gamma(K)\Gamma(\rho+1)\Gamma(w),
\]
using the integer code from
\eqref{eq:integer-descriptor-code}. Assume that this header
is valid with the stated values and that the accompanying array is
in lexicographic order on the complete address set
\[
 [P]\times[2^K]^D\times[S],
 \qquad M=PS\,2^{KD},
\]
with the following record format. For every address
$x\in[P]\times[2^K]^D\times[S]$, write its record as
\[
 z(x)=\sum_{j=0}^{r-1}z_{x,j}y^j,
\]
stored in coefficient order. For every $x$ and $0\le j<r$, assume
$\lvert z_{x,j}\rvert\le1$ and
$\Re z_{x,j},\Im z_{x,j}\in2^{-p}\mathbb Z$.
Each component is $2^{-p}$ times a signed $w$-bit two's-complement
integer, written most significant bit first; the real component
precedes the imaginary component. Require the pointwise bounds
\[
 \log_2(2M)\le C_Mp,\qquad
 |\mathcal E|\le C_{\rm desc}p,\qquad
 p+2\le w\le C_wp,
\]
where $|\mathcal E|$ is the header's bit length. The condition
$w\ge p+2$ allows every permitted disk-grid component, including
$-1$ and $1$. The logical bit volume is $V=2Mrw=\Theta(Mrp)$.

On each of the $D$ consecutive $K$-bit address chunks, select bit
$\rho$, numbered from the least significant bit, and let
$\mathcal H$ be the tensor product of
\[
 H_0=\frac12\begin{pmatrix}1&1\\1&-1\end{pmatrix}
\]
on those selected bits. It acts coefficientwise and fixes every
unselected address bit and both spectator fields. Let $Q_p$ truncate
the real and imaginary parts of every coefficient toward zero to
$2^{-p}\mathbb Z$.

The procedure returns $Q_p(\mathcal H z)$ in the same array,
coefficient order, and component format in $O(Vd^{\lambda'})$ time,
uniformly over all these
choices of $p,d,r,D,P,S,\rho,w$ and $z$. Every returned coefficient belongs
to the complex unit disk, and, in the maximum coefficient norm,
\[
 \|Q_p(\mathcal H z)-\mathcal H z\|_\infty
       <\sqrt2\,2^{-p}.
\]
The procedure evaluates the entire layer exactly before this final
truncation. Put
\[
 C_0=128m\bigl(s_{\rm c}+64(W_{\rm c}+m+1)^3\bigr)^2,
\]
where $m,W_{\rm c},s_{\rm c}$ are the fixed constants of
Proposition~\ref{prop:complex-motif-interface}. The procedure uses
at most $p+\lceil C_0d^{C_1}\rceil$ fractional bits and
$\lceil C_0d^{C_1}\rceil+3$ integer and sign bits per real component
internally; these widths are $O(p)$. The time includes address changes,
deterministic repair, row padding and its removal, and fixed-tape
workspace cleanup. Its tape count is fixed, and its time constant is
independent of the allowed array lengths and coefficient values.
\end{proposition}

\begin{proof}
The procedure reads the current parameters and component width from
$\mathcal E$. Parsing this header and forming the initial collapsed
descriptors take polynomial time in $p$, already charged in the
descriptor work above.
The phase-frame construction reduces a group of $e$ selected bits to
$s_{\rm c}$ children on $e/m$ selected bits.  The row split gives
each child $1/W_{\rm c}$ of the parent's logical volume. The packed
selected-bit procedure supplies the binary basis changes. We verify
that its cost is uniform over the displayed bands for $d,r$.
Put $\Delta=\lceil C_0d^{C_1}\rceil$. Since $d\le b_dp^\epsilon$
and $\epsilon C_1<1$, a common cutoff gives $\Delta\le p$. The guarded
record width is then
$R_{\rm g}=2r(p+2\Delta+3)\le12rp$; also $R_{\rm g}\ge2rp$.
Padding at most doubles the original record count $M$, and row splitting
only decreases it. Thus a packed child with $M_{\rm child}$ records has
the following value $A_*$ of its rectangular width parameter:
\[
\begin{split}
 A_*
 &=\lceil\log_2(2M_{\rm child}R_{\rm g})\rceil\\
 &\le\lceil\log_2(4MR_{\rm g})\rceil\le C_Ap,
 \qquad C_A=C_M+b_r+7.
\end{split}
\]
The last inequality uses $\log_2(2M)\le C_Mp$,
$\log_2r\le b_rp/d\le b_rp$, and $\log_2p\le p$.

Uniformly for large $p$, the bands give $f\le d\le p$ and
\[
 K\ge\tfrac12a_d^c p^{\epsilon c},\qquad
 r\ge2^{(a_r/b_d)p^{1-\epsilon}}.
\]
The first inequality follows from $K=\lfloor d^c\rfloor$ once
$d^c\ge2$. In particular, $K\ge6$ and $K/\log p\to\infty$, while
$r$ dominates every fixed polynomial. Substituting the common bound
$A_*\le C_Ap$ into the two additional normalized
terms in \eqref{eq:packed-rectangle-time} gives
\[
\begin{aligned}
 \frac{A_*^3}{R_{\rm g}}
 &\le\frac{C_A^3p^2}{2r},\\
 80(f-1)2^{-K}A_*\left(1+\frac{A_*}{R_{\rm g}}\right)
 &\le80C_Ap^2 2^{-K}\left(1+\frac{C_A}{2r}\right).
\end{aligned}
\]
Both bounds tend to zero uniformly over all current allowed $d,r$.
Thus Lemma~\ref{lem:packed-selected-bit-rectangle} has the required
$O(M_{\rm child}R_{\rm g}(L^\tau+1))$ cost for every packed call,
with one common time constant.
The row construction and the recurrence analysis
above give total time $O(Vd^{\lambda'})$ on fixed tapes, including
all padded and auxiliary streams. Equation~\eqref{eq:phase-butterfly-layer}
converts the resulting exact $C$-layer to $\mathcal H$.

The guard-width calculation above bounds both denominator and magnitude
growth by $\lceil C_0d^{C_1}\rceil$. Since $d\le b_dp^\epsilon$ and
$\epsilon C_1<1$, this is $o(p)$ uniformly over the allowed $d$, so
every intermediate coefficient is represented exactly with
the stated $O(p)$ width. The initial scan removes or adds only
redundant leading sign bits and appends the reserved fractional zeros
to put each supplied component in this guarded format. To truncate
toward zero at the end, shift its absolute numerator right by the
number of added fractional bits, restore its sign, and write the
result in $w$ bits. These scans cost $O(V)$.
Finally, every row of $\mathcal H$ has
absolute row sum one. Thus $\mathcal H$ maps disk arrays to disk
arrays, and componentwise truncation preserves that disk and has
error less than $\sqrt2\,2^{-p}$: truncation cannot increase the
absolute value of either component, and changes each by less than
$2^{-p}$. The truncated numerators have magnitude at most $2^p$,
so the supplied component width can hold them.
\end{proof}
